\documentclass{article}

\usepackage{mathtools}
\usepackage{geometry}
\usepackage{amssymb}
%Custom Commands
\geometry{
	a4paper,
	total={170mm,257mm},
 	left=20mm,
 	top=20mm
 }

 \mathtoolsset{showonlyrefs}
\begin{document}
	\begin{center}
		\large
		\textbf{Tutorial Sheet 1 Solutions}\\
		\large
		\textbf{ECE F434: Digital Signal Processing}\\
		\normalsize
		\textbf{07/08/2025}
	\end{center}
	\begin{enumerate}
		\item Answer the following Questions 
		\begin{enumerate}
			\item Check the following functions for linearity. 
			\begin{enumerate}
				\item $y[n] = Ax[n] + B$ where $A$ and $B$ are constants.
			
				\textbf{Answer:} Use the linearity property $y[a + b] = y[a] + y[b]$.
					\begin{align}
						\text{Let: }y_1[n] &= Ax_1[n] + B \text{ and } y_2[n] = Ax_2[n] + B \\
						\textbf{Additivity Test: }\\
						y[n] &= A(x_1[n] + x_2[n]) + B \\
						y[n] &= Ax_1[n] + B + Ax_2[n] + B = A(x_1[n] + x_2[n]) + 2B
					\end{align}
					You'll see that the two equations are not equal, as such this function is non linear.
				\item $y[n] = x[n^2]$\\
					\textbf{Answer:} Similarly to above, we see that:
					\begin{align}
						\text{Let: }y_1[n] &= Ax_1[n] + B \text{ and } y_2[n] = Ax_2[n] + B \\
						\textbf{Additivity Test: }\\
						y[n] &= y_1[n] + y_2[n]\\
						y[n] &= x_1[n^2] + x_2[n^2]
					\end{align}
					Given that both of the equations are the same, we can conclude that this function is linear. 	
			\end{enumerate}
		\end{enumerate}
		\item  The following two signals are available. Find the cross correlation.
		\begin{align}
			x[n] = \{... , 0, 0, 2, -1, 3, 7,\mathbf{1} , 2, -3, 0, 0, ...\}\\
			y[n] = \{... , 0, 0, 1, -1, 2, -2, \mathbf{4}, 1, -2, 5, 0, 0, ...\}
		\end{align}		
		\textbf{Answer: } \\
		Cross correlation is given by
		\begin{align}
			r_{xy} = \sum_{n = -\infty}^{\infty} x[n] y[n-l] \forall l
		\end{align}
		As such, for the values of $l$ as 0, 1 and -1:
		\begin{align}
			r_{xy}(0) = 2 + 1 + 6 - 14 + 4 + 2 + 6 = 7 \\
			r_{xy}(1) =-1 - 3 + 14 - 2 + 8 - 3 = 13 \\
			r_{xy}(-1) = -2 - 2 - 6 + 28 + 1 -4 - 14 = 0
		\end{align}
		You may follow the same method for the remaineder of the values of $l$.
		\item The fibonnaci sequence is defined by:
			\begin{align}
				f[n] = f[n-1] + f[n-2] \hspace{1cm}
			\end{align}
			\begin{enumerate}
				\item Develop a direct formula to calculate $f[n]$ only based on $n$
					\textbf{Answer: }\\
					Let $f[n] = \alpha r^n$
					\begin{align}
						\therefore \alpha r^n - \alpha r^{n-1} - \alpha r^{n-2} = 0\\
						\frac{1 \pm \sqrt{5}}{2} \\
						\therefore f[n] = \frac{1}{\sqrt{5}} \left(\frac{1 + \sqrt{5}}{2}\right) - \frac{1}{\sqrt{5}} \left(\frac{1 - \sqrt{5}}{2}\right)
					\end{align}
				\item Show that $y[n]$ can also be represented by the LTI system $y[n] = y[n-1] + y[n-2] + x[n-1]$.

					\textbf{Answer :}\\ 
					When $x[n]$ is an impulse ($\delta (n)$), the sequence begins with 0 at origin, 1 at $n = 1$ and then continues as the fibonnaci seqnuence must go by. As such, this system is  avalid representation of the Fibonnaci series.
			\end{enumerate}
		\item Let $x[n]$ and $y[n]$ be two independent zero mean WSS random signals with autocorrelations $R_{xx}$ and $R_{yy}$. Consider a random singal $v[n]$ obtained by a linear combination of $x$ and $y$. Express that autocorrelation and cross correlations $R_{vx}$, $R_{vy}$ and $R_{vv}$ in terms of $R_{xx}$ and $R_{yy}$.

			\textbf{Answer: }\\

			Let: 
			\begin{equation}
				v[n] = ax[n] + by[n]
			\end{equation}
			From here, we can easily find the autocorrelation:
			\begin{equation}
				R_{vv} = E(v[n+m]v[m]) = a^2R_{xx} + b^2R_{yy}
			\end{equation}
			
	\end{enumerate}
	
\end{document}

